\subsection{Research Objectives}

To advance rigorous ESG ABSA research in multilingual Indonesian/English disclosures and to align with the research questions and pipeline design articulated in previous sections, the following five objectives are stated in precise, measurable, and action-oriented terms. Each objective directly maps to a methodological component of the intended ABSA pipeline and evidences a clear means of evaluation.

% \subsubsection{Objective 1 – Build the OCR-to-record transformation pipeline (PDF-to-structured-ESG representation)}

% To design and implement an end-to-end OCR-enabled pipeline that converts Indonesian and English sustainability reports and regulatory filings (PDF, scanned or native) into a sentence- and paragraph-level, regulator-aligned structured representation (JSON/NER-ready schema) with explicit mappings to GRI/SASB/TCFD concepts.
% To evaluate the pipeline on a diverse set of 60–100 document samples (across Indonesian and cross-language variants) in terms of (a) OCR accuracy, (b) layout-aware extraction quality, (c) fidelity of section and sentence segmentation, and (d) correctness of ESG-topic alignment, reporting regression testing, and cross-language alignment.
% Deliverables: a reproducible OCR-to-record toolchain, annotated preprocessing logs, and a documented data schema linking document regions to ESG aspects, with verifiable performance metrics (character error rate, sentence boundary precision/recall, and taxonomy-mapping accuracy).
% Supporting basis: prior work emphasizes the necessity of robust document structure extraction, OCR/layout-aware processing, and cross-language alignment to support downstream ABSA tasks in multilingual ESG contexts \cite{10.1109/icdar.2011.302,10.1093/jamiaopen/ooac045,10.21203/rs.3.rs-6644838/v1,10.3390/app122111063,10.1108/jd-03-2023-0055,10.1145/3151509.3151525,10.1145/3167132.3167236,10.1145/3322905.3322909}. Climate- and ESG-specific ABSA literature demonstrates the value of anchoring textual content to regulatory concept mappings for interpretability and auditability \cite{10.21203/rs.3.rs-3600821/v1,10.2139/ssrn.3796152,10.1111/1911-3846.12832,10.53555/kuey.v30i5.4019,10.1371/journal.pone.0288052,10.1108/dta-01-2024-0065}

% \subsubsection{Objective 2 – Implement tone-aware ABSA for bilingual ESG disclosures}

% To design and implement a sentence- and aspect-level ABSA pipeline that yields: (i) aspect-term extraction and mapping to ESG pillars (environment, social, governance), (ii) sentence-level sentiment polarity per aspect, and (iii) document-level tonal context to situate per-sentence signals.
% To evaluate the ABSA system using a bilingual Indonesian/English corpus with gold-standard sentence-level annotations, computing precision, recall, F1 for aspect-term extraction, sentiment accuracy per aspect, and calibration metrics (confidence intervals, reliability diagrams) across languages.
% Deliverables: a fully documented ABSA model with domain-adapted fine-tuning (e.g., Indonesian-English multilingual transformer with domain tokens), evaluation report including language-wise performance, and an open-schema for aspect-to-regulatory mapping.
% Rationale: ABSA literature consistently argues for sentence- and aspect-level granularity to achieve actionable ESG signals, and cross-lingual ABSA work underscores the need for robust multilingual mappings and domain adaptation when handling Indonesian/English content \cite{10.1002/csr.70089,10.1016/j.inffus.2025.103073,10.1111/eufm.12509,10.48550/arxiv.2511.00024,10.3390/math12193119}. Tone-aware design advances beyond generic sentiment by ensuring the polarity is tied to concrete ESG aspects, thereby improving interpretability for investors and regulators \cite{10.1108/mf-01-2023-0020,10.2139/ssrn.3738629} .

% \subsubsection{Objective 3 – Compare tone outputs with ClimateBERT-label classifications for climate-related disclosures}

% To conduct a controlled comparison between (a) tone-based ABSA outputs (document/section-level polarity and aspect-level sentiment) and (b) dedicated ClimateBERT-style climate-target labels (net-zero/reduction targets and climate-specific classifications) on identical climate-related passages.
% To quantify concordance/discordance between the two labeling regimes, assess predictive validity for reported climate targets and commitments, and identify systematic divergences that inform model selection or ensemble strategies.
% Deliverables: a comparative study report with statistical analyses (e.g., Cohen’s kappa for agreement, confusion matrices for target labels vs sentiment, and regression analyses linking labels to disclosed targets), plus recommendations for when to fuse or separate climate-target signals from tone signals.
% Rationale: emerging climate-disclosure NLP work demonstrates that domain-specific models (e.g., ClimateBERT variants) capture climate signals beyond generic tone, and comparing these with tone-centric ABSA outputs reveals when specialized targets improve decision-useful signal extraction (\cite{10.21203/rs.3.rs-3600821/v1,10.48550/arxiv.2510.01222,10.18653/v1/2023.emnlp-main.975,10.1108/dta-01-2024-0065}. The objective directly enables evidence-based taxonomy and model selection for ESG analytics in climate-relevant domains.

% \subsubsection{Objective 4 – Evaluate extraction quality and identify failure patterns (diagnostics for ABSA)}

% To develop and deploy a diagnostics framework that automatically identifies extraction weaknesses, including (i) incorrect/omitted ESG-aspect mappings, (ii) misattribution of sentiment polarity to wrong aspects, (iii) language-specific lexical drift and regulatory lexicon gaps, and (iv) failures due to near-duplicate or multi-aspect sentences.
% To quantify failure modes with interpretable metrics (per-sentence error taxonomy, per-document error rates, calibration of confidence scores) and implement a closed-loop process for model improvement (lexicon updates, annotation guideline revisions, re-training or fine-tuning schedules).
% Deliverables: an ABSA failure-analytics dashboard, a formal error taxonomy document, and a protocol for iterative model refinement with traceability to ground-truth annotations and regulatory mappings.
% Rationale: Section 1.3 highlighted the critical need for extraction diagnostics and failure analysis to drive robust ABSA in ESG contexts \cite{10.1002/csr.70089,10.2139/ssrn.3738629,10.1108/mf-01-2023-0020}. A formal diagnostics framework operationalizes this need, enabling continuous improvement and regulator-friendly transparency.

% \subsubsection{Objective 5 – Create reproducible visualization and documentation artifacts for auditability}

% To establish standardized visualization and documentation artifacts that support reproducibility and regulatory auditability, including: (a) an annotated corpus with language tags, aspect taxonomies, and ABSA labels, (b) model configuration records (architecture, hyperparameters, seeds), (c) prompts/templates and in-context example sets for prompts-based LLM components, (d) evaluation results with language- and framework-specific breakdowns, and (e) executable notebooks and data-access guidelines that enable independent replication.
% To evaluate the usefulness of artifacts through a replication study with independent researchers attempting to reproduce key results within predefined time constraints, reporting any deviations and remediation steps.
% Deliverables: a reproducibility package comprising data, code, prompts/templates, schemas, and visualization dashboards, accompanied by a formal reproducibility report validating the ability to reproduce core findings.
% Rationale: Literature on climate- and ESG-NLP emphasizes the importance of transparent methodology, standardized documentation, and reproducible pipelines to ensure trust, accountability, and regulatory auditability in ESG analytics \cite{10.2139/ssrn.3796152,10.1371/journal.pone.0288052,10.1051/e3sconf/202343602004,10.1108/medar-11-2020-1097,10.1108/aaaj-06-2021-5330,10.1108/aaaj-02-2020-4445}. This objective ensures that technical contributions are accessible, auditable, and citable.

% \subsubsection{Objective 6 – Assess cross-model and cross-prompt stability; develop robust ensemble/verification strategies}

% To quantify ABSA output stability across multiple models (monolingual and multilingual transformers, climate-focused variants, general ABSA baselines) and across diverse prompt templates (including zero-shot and few-shot configurations), measuring per-sentence agreement, label stability, and variability in aspect extraction.
% To design and evaluate ensemble strategies (e.g., majority voting, self-consistency prompts, retrieval-augmented generation with external regulatory lexicons) and verification layers that improve reliability and regulator-friendly interpretability of outputs.
% Deliverables: a stability benchmark suite with defined datasets, a set of recommended ensemble/verification configurations, and guidance for practitioners on achieving stable, auditable ESG ABSA outputs in bilingual Indonesian/English reporting.
% Rationale: The need to address prompt-induced instability and model variability in LLM-augmented ABSA is well documented in recent work on ESG NLP and ABSA robustness \cite{10.21203/rs.3.rs-6614220/v1,10.1142/9789811250903_0004,10.30574/ijsra.2024.11.1.0305,10.1609/aaai.v38i21.30574,10.48550/arxiv.2510.01222,10.18653/v1/2023.emnlp-main.975,10.30519/ahtr.1436175,10.1609/aaai.v38i16.29726}. This objective provides formal measurement and practical remedies to ensure reliable deployment in regulatory and investment contexts.

Cross-cutting notes

Language reach: All objectives are designed to operate in Indonesian and English, with explicit mechanisms for cross-language alignment, regulatory lexicon integration (GRI/SASB/TCFD), and cross-framework comparability.
Regulatory alignment: The pipeline outputs will be mappable to OJK/IDX expectations where applicable, ensuring that ABSA signals can support compliance monitoring and investor communications in Indonesia.
Evaluative rigor: Each objective includes concrete metrics, target sample sizes, and deliverables to enable rigorous assessment, replication, and extension in future work.